\chapter{Typed Morphisms, Spans, and Polyxan Rewriting}
\label{ch:typed-morphisms}

This chapter develops the categorical machinery necessary for Polyxan
rewriting theory.  In analogy with how the RSVP formalism (Part~I)
built its dynamics from variational principles and jet prolongations,
we now build the discrete rewrite calculus from typed morphisms, spans,
and curvature-decreasing transformations.

Typed morphisms enforce semantic coherence; spans encode rewriting rules;
the extraction–hoisting–reduction (EHR) cycle defines the canonical
evaluation operator; and curvature monotonicity guarantees termination,
yielding a unique normal form (the media quine) for all finite Polyxan
hypergraphs.

% ============================================================
\section{Typed Morphisms Revisited}
% ============================================================

Let $\mathbf{TypHyper}$ denote the category of typed hypergraphs introduced
in Chapter~\ref{ch:typed-hypergraphs}.  We recall that morphisms
\[
f: \mathcal{G} \to \mathcal{H}
\]
preserve incidence, types, tags, and weights.

Typed morphisms serve three central purposes in this chapter:
\begin{enumerate}
\item they form the vertical arrows in the PolyxSpan double category,
\item they define matches of rewriting rules inside host hypergraphs,
\item they preserve curvature under admissible reduction.
\end{enumerate}

\begin{proposition}[Typed Morphisms Preserve Primitive Curvature]
For any $f : \mathcal{G} \to \mathcal{H}$,
\[
\mathcal{H}_0[\mathcal{G}] \ge \mathcal{H}_0[\mathcal{H}]|_{\mathrm{im}(f)}.
\]
\end{proposition}

Thus morphisms are entropy-non-increasing in the discrete sense.

% ============================================================
\section{Spans of Typed Hypergraphs}
% ============================================================

To rewrite a subgraph of $\mathcal{G}$, we need a way to express
a left-hand pattern $L$ embedded into $\mathcal{G}$ and replaced by
a right-hand pattern $R$.

\begin{definition}[Span]
A \emph{span} in $\mathbf{TypHyper}$ is a diagram
\[
L \xleftarrow{\ell} K \xrightarrow{r} R
\]
where $\ell$ and $r$ are typed monos (injective-on-structure morphisms).
\end{definition}

Here $K$ is the \emph{interface hypergraph} encoding the part of $L$
that is preserved in $R$.

\begin{definition}[Match]
A \emph{match} of $L$ in $\mathcal{G}$ is a typed mono
\[
m : L \hookrightarrow \mathcal{G}.
\]
\end{definition}

Spans are the discrete analogue of local coordinate charts for continuous
rewriting of RSVP fields.  The interface $K$ plays the rôle of the
overlap domain.

% ============================================================
\section{The Double Category \texorpdfstring{$\mathbf{PolyxSpan}$}{PolyxSpan}}
% ============================================================

\begin{definition}[The Double Category $\mathbf{PolyxSpan}$]
Objects: typed hypergraphs.

Vertical morphisms: typed morphisms $f : \mathcal{G}\to\mathcal{H}$.

Horizontal morphisms: spans $L \xleftarrow{\ell} K \xrightarrow{r} R$.

2-cells: commuting squares
\[
\begin{tikzcd}
L \arrow[d,"f_L"'] & K \arrow[l,"\ell"'] \arrow[r,"r"] \arrow[d,"f_K"] & R \arrow[d,"f_R"]\\
L' & K' \arrow[l,"\ell'"'] \arrow[r,"r'"] & R'
\end{tikzcd}
\]
in $\mathbf{TypHyper}$.
\end{definition}

\begin{theorem}
$\mathbf{PolyxSpan}$ is a strict double category with pullback-based
horizontal composition.
\end{theorem}

This structure supports typing constraints, semantic tagging, curvature
control, and eventually the EHR cycle as a horizontal endofunctor.

% ============================================================
\section{Typed Rewriting via Pushouts}
% ============================================================

A rewriting rule is a span $L \xleftarrow{\ell} K \xrightarrow{r} R$.
Given a match $m : L\hookrightarrow\mathcal{G}$, we form the pushout:

\[
\begin{tikzcd}
K \arrow[r,"r"] \arrow[d,"\ell"'] & R \arrow[d]\\
L \arrow[r,"m"] & \mathcal{G}
\end{tikzcd}
\qquad\leadsto\qquad
\begin{tikzcd}
L \arrow[d,"m"'] & R \arrow[l,"\ast"'] \arrow[d]\\
\mathcal{G} & \mathcal{G}' \arrow[l,"\ast"']
\end{tikzcd}
\]

\begin{definition}[Typed Rewrite Step]
A typed rewrite step
\[
\mathcal{G} \;\Rightarrow\; \mathcal{G}'
\]
is the replacement of $m(L)$ in $\mathcal{G}$ by $R$ through a pushout over $K$.
\end{definition}

Typedness of $\ell,r,m$ guarantees:
\[
\mathrm{type},\mathrm{tag},w \quad\text{are preserved or refined in the rewrite}.
\]

% ============================================================
\section{Confluence and the Sesquicategory of Rewriting}
% ============================================================

Let $\mathbf{Rew}(\mathcal{R})$ denote the sesquicategory whose objects are
typed hypergraphs, whose 1-cells are rewrite sequences, and whose
2-cells are commutations of overlapping rewrite steps.

\begin{definition}[Local Confluence]
A rewriting system $\mathcal{R}$ is locally confluent if whenever
\[
\mathcal{G} \Rightarrow \mathcal{H}_1,
\qquad
\mathcal{G} \Rightarrow \mathcal{H}_2,
\]
there exists $\mathcal{K}$ such that
\[
\mathcal{H}_1 \Rightarrow^\ast \mathcal{K},
\qquad
\mathcal{H}_2 \Rightarrow^\ast \mathcal{K}.
\]
\end{definition}

\begin{lemma}[Newman]
If $\mathcal{R}$ is terminating (no infinite rewrite chains), local confluence
implies global confluence.
\end{lemma}

We will prove termination via curvature monotonicity.

% ============================================================
\section{The Extraction–Hoisting–Reduction (EHR) Cycle}
% ============================================================

The EHR cycle is the Polyxan analogue of variational descent:

\begin{itemize}
\item \emph{Extraction}: isolate minimal subgraphs violating type or tag invariants.
\item \emph{Hoisting}: lift structure to higher-typed or canonical forms.
\item \emph{Reduction}: eliminate curvature or tag-entropy defects.
\end{itemize}

\begin{definition}[EHR Endofunctor]
The EHR cycle defines an endofunctor
\[
\mathbb{EHR} : \mathbf{TypHyper} \to \mathbf{TypHyper}.
\]
\end{definition}

\begin{proposition}
$\mathbb{EHR}$ is curvature-decreasing:
\[
\mathcal{H}_0[\mathbb{EHR}(\mathcal{G})]
\le
\mathcal{H}_0[\mathcal{G}],
\]
with strict inequality unless $\mathcal{G}$ is a media quine.
\end{proposition}

This parallels the RSVP Lyapunov functional of Chapter~\ref{ch:stability}.

% ============================================================
\section{Termination via Curvature Monotonicity}
% ============================================================

\begin{theorem}[Termination]
Every sequence of EHR rewrites
\[
\mathcal{G}_0 \Rightarrow \mathcal{G}_1 \Rightarrow \cdots
\]
terminates in finitely many steps.
\]
\end{theorem}

\begin{proof}
$\mathcal{H}_0[\mathcal{G}]$ is a non-negative integer-valued functional and
strictly decreases along nontrivial rewrite steps; thus no infinite descent is
possible.
\end{proof}

Thus Polyxan rewriting enjoys the same global well-posedness enjoyed by RSVP flow.

% ============================================================
\section{Polyxan Energy and Critical Points}
% ============================================================

Define the Polyxan energy
\[
\mathcal{E}[\mathcal{G}]
=
\sum_{e\in E} \delta(e)^2
\;+\;
\sum_{v\in V} H_{\mathrm{tag}}(v)^2.
\]

\begin{proposition}
Media quines are precisely the critical points of $\mathcal{E}$.
\end{proposition}

Thus Polyxan energy plays the rôle of a discrete semantic action.

% ============================================================
\section{Unique Normal Forms}
% ============================================================

We now arrive at the central theorem of Polyxan rewriting.

\begin{theorem}[Polyxan Unique Normal Form Theorem]
Every finite typed hypergraph $\mathcal{G}$ admits a unique normal form
$\mathcal{G}^\ast$, and $\mathcal{G}^\ast$ is a media quine.
\end{theorem}

\begin{proof}
Termination (previous section) plus confluence (Newman lemma) yields a unique
normal form.  The curvature-critical characterization of media quines shows
$\mathcal{G}^\ast$ must be a fixed point of $\mathbb{EHR}$, hence a quine.
\end{proof}

This theorem is the discrete algebraic analogue of the Semantic Relaxation
Theorem in Part~I.

% ============================================================
\section{Summary}
% ============================================================

In this chapter we:
\begin{itemize}
\item constructed the double category $\mathbf{PolyxSpan}$ of typed spans and morphisms,
\item defined typed rewriting rules via pushout replacement,
\item developed the sesquicategory $\mathbf{Rew}(\mathcal{R})$ and proved confluence under termination,
\item introduced the extraction–hoisting–reduction cycle as a curvature-decreasing endofunctor,
\item proved global termination via curvature monotonicity,
\item defined the Polyxan energy as a discrete semantic action,
\item proved the Polyxan Unique Normal Form Theorem.
\end{itemize}

Typed spans and curvature-decreasing rewriting now form the algebraic heart
of Polyxan media.  The next chapter (Chapter~\ref{ch:spectral-semantics})
will introduce spectral invariants and curvature measures to deepen this structure.
